Ambientes de simulação aeroespacial em larga escala geram volumes substanciais de dados
brutos cujo pós-processamento representa um gargalo crítico nos fluxos de trabalho analíticos.
Quando lotes de simulação são processados sequencialmente em uma única máquina, o tempo necessário
para derivar métricas significativas e valores agregados cresce proporcionalmente ao número
de execuções, limitando a agilidade analítica necessária para o Planejamento Baseado em
Capacidades em contextos de defesa.

Este trabalho propõe, implementa e avalia uma infraestrutura de computação distribuída projetada
para acelerar o estágio de pós-processamento de simulações aeroespaciais baseadas em cenários,
produzidas pelo projeto ASA, desenvolvido no Instituto de Estudos Avançados da Força Aérea
Brasileira. O projeto, intitulado AeroSOuDA, combina o Apache Spark (via PySpark) para extração
paralela de métricas com o Hadoop Distributed File System (HDFS) para armazenamento escalável
com localidade de dados, expondo uma interface de linha de comando interativa por meio da qual
analistas definem funções de métrica e de agregação antes do envio do trabalho, sem exigir
conhecimentos computacionais elevados.

Dois modos de implantação são apresentados: um cluster local via Docker Compose para
prototipagem, e uma implantação em máquina virtual remota acessada por um cliente
Python via SSH e SFTP. A avaliação experimental em dois cenários analíticos sobre dois
conjuntos de dados independentes (1,3\,GB e 1,5\,GB, respectivamente) demonstra speedups de
$2{,}86\times$ e $3{,}14\times$ com quatro executores em um cluster de host único,
correspondendo a eficiências paralelas de $71{,}4\%$ e $78{,}6\%$. Um experimento de
escalabilidade com um, dois e quatro executores no conjunto maior confirma que a arquitetura
escala beneficamente, atingindo $93{,}4\%$ de eficiência com dois executores. O overhead de
comunicação SSH permanece inferior a $1\%$ do tempo total de execução em ambas as campanhas,
confirmando a viabilidade do modelo de implantação remota.

Os resultados estabelecem que o paradigma MapReduce, implementado via PySpark e HDFS,
fornece uma base prática e escalável para o pós-processamento distribuído de lotes de
simulação em um contexto aeroespacial de defesa, reduzindo o tempo de execução em
$65$--$68\%$ em relação ao referencial sequencial, mantendo o conjunto de dados bruto
integralmente no servidor. O speedup medido cresce com o tamanho do conjunto de dados,
consistente com as previsões da Lei de Amdahl derivadas do experimento de escalabilidade.